\section{2025 AI 赛道复盘：Coding、视频、Agentic Commerce 与垂直应用}

本章从公司个案回到赛道。Freda 提出一个投资圈的公开秘密：如果一个领域一年内创造出 10 亿美元年化收入，不管之前理解不理解，都必须认真研究，因为市场在告诉你这里可能产生世界级公司。2025 年最典型的是 Coding，一年内创造超过 50 亿美元年化收入。

\begin{figure}[H]
\centering
\includegraphics[width=0.96\textwidth]{ai-sector-map-2025.png}
\caption{2025 AI 赛道地图：Coding、视频、Agentic Commerce、金融/法律、医疗/客服和 AI 应用估值。自制概念图，依据 01:03:43--01:09:51 对谈内容整理。}
\end{figure}

\begin{knowledgebox}{读图：赛道热度要看收入速度}
Coding 是最清晰的收入爆发赛道；视频因为可以直接优化内容对象，可能改变媒体行业；Agentic Commerce 会影响广告、电商、支付；金融、法律、医疗、客服等垂直应用则看 ROI 和行业数据。
\end{knowledgebox}

\subsection{Coding 为什么最适合模型公司亲自做}

本节先拆 Coding，因为它是 2025 年最清晰的 AI 应用收入样本，也最能说明哪些垂直场景会被模型公司亲自吃掉。

Coding 与金融、法律、医疗不同，它不需要额外采集大量行业知识，模型公司天然拥有代码数据、开发者入口和工具链。Claude Code、OpenAI Codex 等产品发布后收入直线上升，说明模型公司在这个垂直领域有很强优势。但竞争格局仍不清晰：如果 Google 或其他模型公司把 Coding 能力免费化，对 Cursor 等独立公司会形成压力。

\begin{warningbox}{垂直应用不一定都适合创业公司}
上一轮软件常从通用平行层再到垂类；AI 时代通用层被大模型吃掉，很多创业机会转向垂直场景。但 Coding 是例外：它太贴近模型公司的训练数据和开发者生态，独立公司必须找到更强产品体验或工作流壁垒。
\end{warningbox}

\subsection{视频：从匹配函数到可优化对象}

Freda 看好视频，因为媒体行业体量巨大，视频是高带宽输入形式，而且视频第一次可以直接被模型优化。TikTok 时代优化的是匹配：让创作者产出视频，再匹配给用户；生成式视频时代，视频本身可以被目标函数优化，例如参与度、广告点击率、转化率或品牌效果。这个变化会让视频从内容分发问题，变成内容生成与优化问题。

\subsection{Agentic Commerce：广告、电商和支付的重排}

Agentic Commerce 指 Agent 替用户购买商品：理解需求、比较商品、推荐、下单和支付。Freda 认为这对商家可能是噩梦，因为商家希望用户直接到自己网站，从而 cross-sell、卖广告和控制流量。Booking.com 依赖 Google 引流，Amazon 拥有巨大的广告收入，这些都是用户直接访问平台才成立的模式。

\begin{importantbox}{Agentic Commerce 的受益者和受损者}
受益者可能是过去买不起广告的小众商家，因为 Agent 可以精准匹配长尾商品；受损者可能是依赖入口流量、广告位和支付中间环节的平台。支付公司也会被重新定价，因为 Agent 自己可以成为用户的钱包和授权入口。
\end{importantbox}

\subsection{AI 应用公司估值：毛利率还是绝对利润}

本节回到估值口径，因为 AI 应用不像传统软件那样天然拥有稳定高毛利模板。投资人必须重新判断：到底该看百分比毛利率，还是看单用户毛利和利润绝对额。

传统软件像“鸡胸肉”：毛利率和净收入留存率高度一致，投资人能流水线估值。AI 应用则不同：使用越多，推理成本越高，毛利率反而可能下降。Freda 的建议是看利润的绝对金额，因为 AI 合同金额和每用户毛利可能比传统软件大得多。AWS 早期毛利率也不如软件，但市场最终比许多软件更大。

\subsection{本章小结}

2025 AI 赛道复盘的核心是收入速度。Coding 已经验证，视频和 Agentic Commerce 可能打开更大互联网重排，垂直应用则要看行业 ROI。AI 应用估值不能只用旧 SaaS 毛利率模板。

